\section{Contratos e Fluxos de Integração}

Os contratos da NSS 1.0 deixam de ser apenas recomendações abstratas. As versões normativas em Markdown estão em \texttt{contracts/api-v1.md}, \texttt{contracts/auth-v1.md} e \texttt{contracts/serving-v1.md}.

\subsection{Frontend--Java}

O frontend consome exclusivamente a API Java sob o prefixo \texttt{/api/v1}. Para a V1.0, o contrato cobre:

\begin{itemize}
    \item login, refresh e logout;
    \item lista de doenças disponíveis;
    \item agregação por município;
    \item agregação por distrito;
    \item agregação por bairro/localidade;
    \item filtros de doença, ano, mês, sexo e faixa etária;
    \item semântica explícita de cobertura disponível/indisponível.
\end{itemize}

A API deve utilizar um campo de métrica semântico, permitindo que a UI mostre \textbf{Notificações} enquanto não houver regra aprovada para \textbf{Casos confirmados}.

\subsection{Autenticação}

A V1.0 utiliza login por e-mail/senha, JWT de curta duração e refresh token opaco rotacionável em cookie HttpOnly/Secure. O frontend mantém o access token em memória e usa \texttt{Authorization: Bearer}. Auto-registro público fica desabilitado.

\subsection{Pipeline--PostgreSQL}

O schema \texttt{analytics} e a Gold Serving V1 funcionam como contratos internos. Mudanças requerem migration Alembic e validação das queries do Java.

\subsection{Java--PostgreSQL}

O Java não precisa reconstruir lógica SINAN. Ele recebe dados Serving e executa agregações relacionais. Regras como codificação de idade, CNES e classificação epidemiológica pertencem à pipeline/contrato de dados, não ao controller.

\subsection{Fluxo de consulta autenticada}

\begin{lstlisting}[language=Mermaid]
sequenceDiagram
    participant U as Usuario
    participant F as Frontend
    participant C as Caddy
    participant J as Java
    participant D as PostgreSQL

    U->>F: login email/senha
    F->>C: POST /api/v1/auth/login
    C->>J: proxy
    J-->>F: access JWT + refresh cookie HttpOnly
    F->>J: GET /api/v1/epidemiology/... + Bearer
    J->>D: SUM/GROUP BY
    D-->>J: agregados
    J-->>F: JSON V1
    F-->>U: mapa/painel
\end{lstlisting}

\subsection{Fluxo de atualização}

\begin{lstlisting}[language=Mermaid]
sequenceDiagram
    participant S as SINAN/CNES
    participant P as Pipeline
    participant D as PostgreSQL
    participant G as Grafana

    P->>S: obter fontes
    S-->>P: dados
    P->>P: Bronze / Silver / Gold Ad Hoc / Serving
    P->>P: validar contrato e reconciliar
    P->>D: publicar Serving V1
    D-->>P: resultado da transacao
    G->>D: consultar metricas de auditoria
\end{lstlisting}
